\documentclass[12pt,a4paper]{article}

\usepackage[margin=2cm]{geometry}
\usepackage{amsmath,amssymb,mathtools}
\usepackage{physics} % for \ket, \bra, \outerproduct, etc.
\usepackage{bm}
\usepackage{cite}
\usepackage{hyperref}

\title{Parity selection by spin-$z$ jump operators acting on \texorpdfstring{$\ket{++}$}{|++>}}
\author{}
\date{}

\begin{document}
	\maketitle
	
	\begin{abstract}
		We consider a two-qubit system subjected to quantum jumps generated by the two collapseing operators
		\(
		\hat{C}_{\pm}\propto S_z^{(1)}\pm S_z^{(2)}
		\),
		with the physical convention
		\(\sigma_z\ket{1}=+\ket{1}\), \(\sigma_z\ket{0}=-\ket{0}\).
		Starting from the product state \(\ket{++}\) (the $+1$ eigenstate of \(\sigma_x\) on each qubit), we show that the first jump projects the system onto a definite $z$-parity subspace.
		A ``positive'' click \((C_+)\) sends \(\ket{++}\) to \(-\ket{\phi_-}\) and then cycles deterministically
		\(-\ket{\phi_-}\to \ket{\phi_+}\to -\ket{\phi_-}\to\ldots\). 
		A ``negative'' click \((C_-)\) sends \(\ket{++}\) to \(-\ket{\psi_-}\) and then cycles
		\(-\ket{\psi_-}\to \ket{\psi_+}\to -\ket{\psi_-}\to\ldots\).
		This selection rule follows from the explicit commutation \([C_{\pm},\Pi_z]=0\) with the $z$-parity operator \(\Pi_z=\sigma_z^{(1)}\otimes\sigma_z^{(2)}\).
		We also derive the jump probabilities at each step. Finally we analyze both concurrence $\mathcal C$ and spin squeezing parameter $\xi^2$ and derive ensemble averages.
	\end{abstract}
	
	\section{Conventions, basis and Bell states}
	We work with two qubits in the computational $z$-basis
	\(\{\ket{1}\equiv\ket{\uparrow_z},\ \ket{0}\equiv\ket{\downarrow_z}\}\),
	with Pauli matrices \(\sigma_\alpha^{(i)}\) acting on qubit \(i=1,2\) and
	\(\sigma_z\ket{1}=+\ket{1}\), \(\sigma_z\ket{0}=-\ket{0}\).
	Spin-$1/2$ operators are \(S_\alpha^{(i)}=\tfrac{1}{2}\sigma_\alpha^{(i)}\).
	The single-qubit $x$-eigenstates are
	\[
	\ket{+}=\frac{\ket{0}+\ket{1}}{\sqrt{2}},
	\qquad
	\ket{-}=\frac{\ket{0}-\ket{1}}{\sqrt{2}},
	\]
	so the initial state is the product
	\[
	\ket{++}=\ket{+}\otimes\ket{+}
	=\frac{1}{2}\!\left(\ket{00}+\ket{01}+\ket{10}+\ket{11}\right).
	\] \\
	We use the standard Bell states
	\begin{align}
		\ket{\phi_\pm}&=\frac{\ket{00}\pm\ket{11}}{\sqrt{2}}
		\quad(\text{even $z$-parity}),\\
		\ket{\psi_\pm}&=\frac{\ket{01}\pm\ket{10}}{\sqrt{2}}
		\quad(\text{odd $z$-parity}).
	\end{align}
	Define the $z$-parity operator
	\[
	\Pi_z=\sigma_z^{(1)}\otimes\sigma_z^{(2)}.
	\]
	Its eigenvalues are \(+1\) on the even subspace (spin aligned), \(\mathrm{span}\{\ket{00},\ket{11}\}=\mathrm{span}\{\ket{\phi_+},\ket{\phi_-}\}\) and \(-1\) on the odd subspace (spin unalinged), \(\mathrm{span}\{\ket{01},\ket{10}\}=\mathrm{span}\{\ket{\psi_+},\ket{\psi_-}\}\).
	
	\section{Quantum jumps dynamics}
	
	\subsection{Collapse (jump) operators and commutation with parity}\label{sec:comm}
	By adding a 50/50 BS the two collapsing operators are the balanced sum/difference of local $z$-spins with an overall rate factor:
	\begin{equation}
		\hat{C}_\pm=\sqrt{\gamma}\,\frac{S_z^{(1)}\pm S_z^{(2)}}{\sqrt{2}}
		=\sqrt{\gamma}\,\frac{\sigma_z^{(1)}\pm \sigma_z^{(2)}}{2\sqrt{2}}.
		\label{Cpm}
	\end{equation}
	Any global prefactor only affects rates, not the direction of the post-jump state.
	
	Since \(\sigma_z^{(1)}\) and \(\sigma_z^{(2)}\) commute and \((\sigma_z^{(i)})^2=\mathbb{I}\),
	\begin{align}
		\big[\sigma_z^{(1)},\Pi_z\big]
		&=\sigma_z^{(1)}\sigma_z^{(1)}\sigma_z^{(2)}-\sigma_z^{(1)}\sigma_z^{(1)}\sigma_z^{(2)}=0,\\
		\big[\sigma_z^{(2)},\Pi_z\big]
		&=\sigma_z^{(2)}\sigma_z^{(1)}\sigma_z^{(2)}-\sigma_z^{(1)}\sigma_z^{(2)}\sigma_z^{(2)}
		=\sigma_z^{(1)}-\sigma_z^{(1)}=0.
	\end{align}
	Therefore \([\sigma_z^{(1)}\pm \sigma_z^{(2)},\Pi_z]=0\), and by \eqref{Cpm} we have
	\begin{equation}
		[\hat{C}_\pm,\Pi_z]=0.
	\end{equation}
	A jump produced by \(\hat{C}_\pm\) preserves the $z$-parity: if a pre-jump pure state \(\ket{\psi}\) is a parity eigenstate, so is the post-jump normalized state : 
	\(\ket{\psi'}=\hat{C}_\pm\ket{\psi}/\norm{\hat{C}_\pm\ket{\psi}}\).
	
	\subsection{First jump from \texorpdfstring{$\ket{++}$}{|++>}: projection to definite parity subspace}
	For pure-state jumps we use the standard update rule
	\begin{equation}
		%\ket{\psi'}=\frac{C_k\ket{\psi}}{\sqrt{\matrixel{\psi}{C^{\dagger} }_k {\psi}}}\,,
		\qquad p_k=\matrixel{\psi}{C_k^\dagger C_k}{\psi}\,\mathrm{d}t,
		\label{eq:jump_rule}
	\end{equation}
	with \(k\in\{+,-\}\) and \(\mathrm{d}t\) a short time step.
	
	It is convenient to evaluate the \emph{direction} of the unnormalized state \(\hat{C}_\pm\ket{++}\) using the numerators of \eqref{Cpm}:
	\begin{align}
		(\sigma_z^{(1)}+\sigma_z^{(2)})\ket{++}
		&=2\big(\ket{11}-\ket{00}\big)
		=-2\sqrt{2}\,\ket{\phi_-},\label{eq:sigplus_on_pp}\\
		(\sigma_z^{(1)}-\sigma_z^{(2)})\ket{++}
		&=2\big(\ket{10}-\ket{01}\big)
		=-2\sqrt{2}\,\ket{\psi_-}.
		\label{eq:sigminus_on_pp}
	\end{align}
	Hence the \emph{normalized} post-jump states are
	\begin{equation}
		\ket{++}\xrightarrow{C_+}-\ket{\phi_-},
		\qquad
		\ket{++}\xrightarrow{C_-}-\ket{\psi_-}.
		\label{eq:first_step}
	\end{equation}
	Overall minus signs are global phases and carry no physical effect; we keep them to make the subsequent alternation visually clear.
	
	Equations \eqref{eq:first_step} show that the first click already fixes the $z$-parity: $C_+$ lands in the even sector (Bell-$\phi$ states), while $C_-$ lands in the odd sector (Bell-$\psi$ states).
	
	\subsection{Subsequent jumps: two-state cycles within each parity sector}
	Using again only the numerators (because normalization is restored at each jump), we find the action on Bell states:
	\begin{align}
		\hat{C}_+\ket{\phi_-}
		&=-\,\ket{\phi_+},&
		\hat{C}_+\ket{\phi_+}
		&=-\,\ket{\phi_-},\\[2mm]
		\hat{C}_-\ket{\psi_-}
		&=-\,\ket{\psi_+},&
		\hat{C}_-\ket{\psi_+}
		&=-\,\ket{\psi_-}.
	\end{align}
	After renormalization, this yields the deterministic cycles
	\begin{equation}
		-\ket{\phi_-}\xrightarrow{C_+}\ket{\phi_+}\xrightarrow{C_+}-\ket{\phi_-}\xrightarrow{C_+}\cdots,
		\qquad
		-\ket{\psi_-}\xrightarrow{C_-}\ket{\psi_+}\xrightarrow{C_-}-\ket{\psi_-}\xrightarrow{C_-}\cdots.
		\label{eq:cycles}
	\end{equation}
	Because of $[C_\pm,\Pi_z]=0$, the ``wrong'' channel is forbidden inside each parity subspace.
	
	\subsection{Jump probabilities at each step}
	From \eqref{Cpm},
	\begin{equation}
		C_\pm^\dagger C_\pm
		=\frac{\gamma}{2}\big(S_z^{(1)}\pm S_z^{(2)}\big)^2
		=\frac{\gamma}{2}\left[\big(S_z^{(1)}\big)^2+\big(S_z^{(2)}\big)^2\pm2S_z^{(1)}S_z^{(2)}\right].
	\end{equation}
	Using \(S_z^2=\tfrac{1}{4}\,\mathbb{I}\) and \(S_z^{(1)}S_z^{(2)}=\tfrac{1}{4}\,\sigma_z^{(1)}\sigma_z^{(2)}=\tfrac{1}{4}\,\Pi_z\), we obtain the identity:
	\begin{equation}
		C_\pm^\dagger C_\pm
		=\frac{\gamma}{4}\left(\mathbb{I}\pm \Pi_z\right).
		\label{eq:rate_identity}
	\end{equation}
	Therefore, for any pure state \(\ket{\psi}\), the transition probability rate is:
	\begin{equation}
		p_\pm(\ket{\psi})=\matrixel{\psi}{C_\pm^\dagger C_\pm}{\psi}\,\mathrm{d}t
		=\frac{\gamma}{4}\Big(1\pm \ev{\Pi_z}_{\psi}\Big)\,\mathrm{d}t.
		\label{eq:general_rate}
	\end{equation}
	For the state \(\ket{++}\), since \(\ev{\sigma_z}_{\ket{+}}=0\), we have \(\ev{\Pi_z}_{\ket{++}}=\ev{\sigma_z}_{\ket{+}}\,\ev{\sigma_z}_{\ket{+}}=0\), hence
	\begin{equation}
		p_+(\ket{++})=p_-(\ket{++})=\frac{\gamma}{4}\,\mathrm{d}t.
		\label{eq:rates_pp}
	\end{equation}
	The two channels are equiprobable on the initial state, and the total jump rate is \(\tfrac{\gamma}{2}\,\mathrm{d}t\).

	For any even-parity state (eigenvalue \(+1\) of \(\Pi_z\); in particular \(\ket{\phi_\pm}\)),
	\begin{equation}
		p_+(\text{even})=\frac{\gamma}{2}\,\mathrm{d}t,\qquad p_-(\text{even})=0.
	\end{equation}
	For any odd-parity state (eigenvalue \(-1\); in particular \(\ket{\psi_\pm}\)),
	\begin{equation}
		p_-(\text{odd})=\frac{\gamma}{2}\,\mathrm{d}t,\qquad p_+(\text{odd})=0.
	\end{equation}
	Thus, once the first click has selected a parity subspace, only the corresponding channel can ever fire; conditioned on a click, the state update within that sector is deterministic as in \eqref{eq:cycles}.
	
	\subsection{State by state transitions}
	Let us collect both the post-jump states and the \emph{unconditional} jump probabilities over a time step \(\mathrm{d}t\):
	\begin{center}
		\renewcommand{\arraystretch}{1.3}
		\begin{tabular}{c|c|c}
			Pre-jump state & Channel & Post-jump state \& probability \\ \hline
			\(\ket{++}\) & \(C_+\) & \(-\ket{\phi_-}\), \(\;p_+=\dfrac{\gamma}{4}\mathrm{d}t\) \\
			\(\ket{++}\) & \(C_-\) & \(-\ket{\psi_-}\), \(\;p_-=\dfrac{\gamma}{4}\mathrm{d}t\) \\ \hline
			\(-\ket{\phi_-}\) & \(C_+\) & \(\ket{\phi_+}\), \(\;p_+=\dfrac{\gamma}{2}\mathrm{d}t\) \\
			\(\ket{\phi_+}\) & \(C_+\) & \(-\ket{\phi_-}\), \(\;p_+=\dfrac{\gamma}{2}\mathrm{d}t\) \\
			(even sector) & \(C_-\) & forbidden, \(\;p_-=0\) \\ \hline
			\(-\ket{\psi_-}\) & \(C_-\) & \(\ket{\psi_+}\), \(\;p_-=\dfrac{\gamma}{2}\mathrm{d}t\) \\
			\(\ket{\psi_+}\) & \(C_-\) & \(-\ket{\psi_-}\), \(\;p_-=\dfrac{\gamma}{2}\mathrm{d}t\) \\
			(odd sector) & \(C_+\) & forbidden, \(\;p_+=0\)
		\end{tabular}
	\end{center}
	The selection rule is a direct consequence of \([C_\pm,\Pi_z]=0\): jumps cannot mix parity sectors. Equation~\eqref{eq:rate_identity} makes this quantitative.
	
	\subsection{Variances within each channel's cycle}
	
	For Bell states the correlation matrix is diagonal:
	\[
	\expval*{\hat\sigma_\alpha^{(1)}\!\otimes \hat\sigma_\alpha^{(2)}}=
	\begin{cases}
		diag(+1, -1, +1) & \text{for }\ket{\phi_+}\ \text{on }(\alpha=x,y,z),\\
		diag(-1, +1, +1) & \text{for }\ket{\phi_-},\\
		diag(+1, +1, -1) & \text{for }\ket{\psi_+},\\
		diag(-1, -1, -1) & \text{for }\ket{\psi_-}.
	\end{cases}
	\]
	Using
	\[
	\hat J_\alpha^2=\tfrac{1}{2}\Big(\hat{\mathbb{I}}+\hat\sigma_\alpha^{(1)}\!\otimes \hat\sigma_\alpha^{(2)}\Big),
	\qquad \expval*{\hat J_\alpha}=0\ \text{for Bell states},
	\]
	one gets
	\[
	\Delta J_\alpha^2=\expval*{\hat J_\alpha^2}
	=\tfrac{1}{2}\Big(1+\expval*{\hat\sigma_\alpha^{(1)}\!\otimes \hat\sigma_\alpha^{(2)}}\Big).
	\]
	\medskip
	So we summarize in a table of variances \(\Delta J_\alpha^2\), for the two cycles):
	\[
	\begin{array}{c|ccc}
		\text{state} & \Delta J_x^2 & \Delta J_y^2 & \Delta J_z^2\\
		\hline
		\ket{\phi_-} & 0 & 1 & 1\\
		\ket{\phi_+} & 1 & 0 & 1\\
		\hline
		\ket{\psi_-} & 0 & 0 & 0\\
		\ket{\psi_+} & 1 & 1 & 0
	\end{array}
	\]	
	
	
	%%%%%%%%%%%%%%%%%%%%%%%%%%%%%%%%%%%%%%%%%%%%%%%%%%%%%%%%%%%%%%%%%%%%%%%%%%%%%%%%%%%%%%%%%%%%%%%%%%%%
	%%%%%%%%%%%%%%%%%%%%%%%%%%%%%%%%%%%%%%%%%%%%%%%%%%%%%%%%%%%%%%%%%%%%%%%%%%%%%%%%%%%%%%%%%%%%%%%%%%%%
	
	
	
	\section{Ensemble-averaged concurrence over trajectories}
	
	For a pure two-qubit state \(\ket{\psi}=a\ket{00}+b\ket{01}+c\ket{10}+d\ket{11}\),
	the concurrence is \(\mathcal{C}(\ket{\psi})=2\,\lvert ad-bc\rvert\). For a separable state it is 0 and for any Bell states it is 1.
	Hence \(\mathcal{C}(\ket{++})=0\) and \(\mathcal{C}(\ket{\phi_\pm})=\mathcal{C}(\ket{\psi_\pm})=1\).
	
	Along each trajectory the concurrence is a step process: it remains \(0\) until the \emph{first} click, and jumps to \(1\) thereafter (all subsequent clicks keep the state maximally entangled within the same parity subspace).
	
	The waiting time to the first click is exponential with rate
	\[
	\lambda=\Tr\!\Big[\big(\hat C_+^\dagger \hat C_+ + \hat C_-^\dagger \hat C_-\big)\, \cdot \rho_{++}\Big]=\frac{\gamma}{2},
	\]
	so the survival probability of no click up to time \(t\) is \(P_0(t)=e^{-\lambda t}=e^{-(\gamma/2)t}\).
	Hence the ensemble-averaged concurrence over many trajectories is
	\[
	\boxed{\ \overline{\mathcal{C}}(t)= prob (\text{at least one click by }t)=1-e^{-(\gamma/2)t}\ } 
	\]
	With \(\gamma=1\), \(\overline{\mathcal{C}}(t)=1-e^{-t/2}\).
	

	
	\section{Spin squeezing: Kitagawa-Ueda vs Wineland}
	
	Define collective spins (for \(N=2\))
	\[
	\hat J_\alpha=\tfrac{1}{2}\big(\hat\sigma_\alpha^{(1)}+\hat\sigma_\alpha^{(2)}\big),\qquad \alpha\in\{x,y,z\}.
	\]
	
	\paragraph{Kitagawa--Ueda (KU).}
	\[
	\xi_S^2 \equiv \frac{4}{N}\,\min_{\mathbf{n}\perp \expval*{\hat{\mathbf{J}}}}
	\Delta J_{\mathbf{n}}^2,\qquad
	\Delta J_{\mathbf{n}}^2=\expval*{\hat J_{\mathbf{n}}^2}-\expval*{\hat J_{\mathbf{n}}}^2.
	\]
	For the Bell states, \(\expval*{\hat{\mathbf{J}}}=\mathbf{0}\) and at least one orthogonal direction has zero variance. Therefore
	\[
	\boxed{\ \xi_S^2(\ket{\phi_\pm})=\xi_S^2(\ket{\psi_\pm}) = 0} 
	\]
	Before the first click the state is \(\ket{++}\), a CSS state aligned along \(x\), with
	\(\min_{\mathbf{n}\perp \hat x}(\Delta J_{\mathbf{n}}^2)=\dfrac{N}{4}=\dfrac{1}{2}\), hence \(\xi_S^2(\ket{++})=1\) as expected.
	Averaging over many trajectories,
	\[
	\boxed{\ \overline{\xi_S^2}(t)=P_0(t)\cdot 1+(1-P_0(t))\cdot 0=e^{-(\gamma/2)t}\ } 
	\]
	
	\paragraph{Wineland (phase-sensitivity) parameter.}
	\[
	\xi_R^2 \equiv \frac{\min_{\mathbf{n}\perp \expval*{\hat{\mathbf{J}}}}
		\Delta J_{\mathbf{n}}^2}{\norm{\expval*{\hat{\mathbf{J}}}}^2}.
	\]
	All Bell states satisfy \(\norm{\expval*{\hat{\mathbf{J}}}}=0\); thus \(\xi_R^2\) is undefined since it diverges).
	In this parity-selected dynamics the KU definition is the meaningful one, whereas the Wineland parameter cannot be used because the mean spin vanishes.
	%%%%%%%%%%%%%%%%%%%%%%%%%%%%%%%%%%%%%%%%%%%%%%%%%%%%%%%%%%%%%%%%%%%%%%%%%%%%%%%%%%%%%%%%%%%%%%%
	%%%%%%%%%%%%%%%%%%%%%%%%%%%%%%%%%%%%%%%%%%%%%%%%%%%%%%%%%%%%%%%%%%%%%%%%%%%%%%%%%%%%%%%%%%%%%%%	
	
	\section{Adding efficiency \texorpdfstring{$\eta$}{eta}}
	
	\subsection{Observed and unobserved part}
	Assume identical detection efficiencies on the two outputs, \(\eta_1=\eta_2\equiv\eta\in[0,1]\).
	The physical jump operators are unchanged,
	\[
	\hat C_{\pm}=\sqrt{\gamma}\,\frac{\hat S^{(1)}_z\pm \hat S^{(2)}_z}{\sqrt{2}}
	=\sqrt{\gamma}\,\frac{\hat\sigma^{(1)}_z\pm \hat\sigma^{(2)}_z}{2\sqrt{2}},
	\qquad
	\hat\Pi_z=\hat\sigma^{(1)}_z\otimes\hat\sigma^{(2)}_z,
	\]
	and the unconditional master equation (average evolution) is unaffected by \(\eta\).
	At the level of the unravelling, each physical channel splits into an \emph{observed}
	and an \emph{unobserved} part:
	\[
	\text{observed: }\sqrt{\eta}\,\hat C_{\pm},\qquad
	\text{unobserved: }\sqrt{1-\eta}\,\hat C_{\pm}.
	\]
	
	\subsection{Jump probabilities as post-jump norms (observed vs physical)}
	
	Using \(\hat C_{\pm}^\dagger\hat C_{\pm}=\tfrac{\gamma}{4}(\hat{\mathbb{I}}\pm \hat\Pi_z)\),
	the physical and observed click probabilities over \(\mathrm dt\) are
	\[
	p_{\pm}^{\mathrm{phys}}(\psi)=\langle \psi|\hat C_{\pm}^\dagger\hat C_{\pm}|\psi\rangle\,\mathrm dt
	=\frac{\gamma}{4}\big(1\pm \langle \hat\Pi_z\rangle_{\psi}\big)\mathrm dt,
	\qquad
	p_{\pm}^{\mathrm{obs}}(\psi)=\eta\,p_{\pm}^{\mathrm{phys}}(\psi).
	\]
	From \(\ket{++}\) one has \(\langle \hat\Pi_z\rangle=0\) and therefore
	\(p_{+}^{\mathrm{obs}}=p_{-}^{\mathrm{obs}}=\tfrac{\eta\gamma}{4}\,\mathrm dt\), total observed rate \(\eta\gamma/2\).
	Because \([\hat C_{\pm},\hat\Pi_z]=0\), inside the even (odd) subspace the \(-\) (\(+\)) channel remains strictly forbidden even for \(\eta<1\).
	
	
	
	\subsection{Concurrence with finite efficiency}
	
	For a pure two-qubit state \(\ket{\psi}=a\ket{00}+b\ket{01}+c\ket{10}+d\ket{11}\),
	\(\mathcal C(\ket{\psi})=2\,|ad-bc|\);
	thus \(\mathcal C(\ket{++})=0\) and \(\mathcal C(\text{Bell})=1\).
	Along each trajectory the first jump makes \(\mathcal C\) jump from \(0\) to \(1\) and it stays there; averaging over physical trajectories it gives
	\[
	\overline{\mathcal C}_{\mathrm{phys}}(t)=1-e^{-(\gamma/2)t}\qquad(\text{independent of }\eta).
	\]
	For the conditional estimate (based on the observed record), between two observed clicks the state in each subspace is a mixture of two orthogonal Bell states with weights \(p\) and \(1-p\).
	For such a Bell-diagonal form one has
	\[
	\mathcal C_{\mathrm{cond}}(t)=|2p(t)-1|=e^{-(1-\eta)\gamma\,(t-t_0)}.
	\]
	Observed clicks arrive with rate \(\lambda=\eta \cdot \gamma/2\) and reset \(\mathcal C_{\mathrm{cond}}\) to \(1\).
	In stationary conditions the time since the last observed click, \(\tau\), is exponentially distributed with density \(\lambda e^{-\lambda \tau}\);
	therefore
	\[
	\langle \mathcal C_{\mathrm{cond}}\rangle_{\mathrm{ss}}
	=\int_{0}^{\infty} e^{-(1-\eta)\gamma\,\tau}\,\lambda e^{-\lambda \tau}\,\mathrm d\tau
	=\frac{\lambda}{\lambda+(1-\eta)\gamma}
	=\boxed{\ \frac{\eta}{2-\eta}\ }.
	\]
	With \(\gamma=1\), \(\lambda=\eta/2\) and the same result follows.
	
	\section{Spin squeezing with finite efficiency}
	
	Let's discuss how it Kitagawa-Ueda spin squeezing behaves in the finite efficiency case. By using the results before, we show that in the even sector, using \((\Delta J_x^2,\Delta J_y^2,\Delta J_z^2)=(0,1,1)\) for \(\ket{\phi_-}\) and \((1,0,1)\) for \(\ket{\phi_+}\),
	the mixture \(\rho_{\mathrm{even}}(t)=p\ket{\phi_{\mathrm{start}}}\!\bra{\phi_{\mathrm{start}}}+(1-p)\ket{\phi_{\mathrm{other}}}\!\bra{\phi_{\mathrm{other}}}\) has
	\[
	\Delta J_x^2=p,\qquad \Delta J_y^2=1-p,\qquad \Delta J_z^2=1,
	\]
	so that \(\min_{\mathbf n}\Delta J_{\mathbf n}^2=\min\{p,1-p\}=\tfrac{1}{2}\big(1-|2p-1|\big)\).
	Therefore
	\[
	\xi_{S,\mathrm{even}}^2(t)=1-|2p(t)-1|=1-e^{-(1-\eta)\gamma\,(t-t_0)}.
	\]
	Every observed \(+\) click resets \(\xi_{S,\mathrm{even}}^2\) to \(0\); between clicks it increases towards \(1\).
	Averaging over the exponential distribution of \(\tau\) with rate \(\lambda=\eta\gamma/2\) yields
	\[
	\big\langle \xi_{S,\mathrm{even}}^2\big\rangle_{\mathrm{ss}}
	=\int_{0}^{\infty}\big(1-e^{-(1-\eta)\gamma\,\tau}\big)\,\lambda e^{-\lambda \tau}\,\mathrm d\tau
	=\boxed{\ 1-\frac{\eta}{2-\eta}=\frac{2(1-\eta)}{2-\eta}\ }.
	\]
	
	In the odd sector, using \((\Delta J_x^2,\Delta J_y^2,\Delta J_z^2)=(0,0,0)\) for \(\ket{\psi_-}\) and \((1,1,0)\) for \(\ket{\psi_+}\),
	any mixture \(p\ket{\psi_-}\!\bra{\psi_-}+(1-p)\ket{\psi_+}\!\bra{\psi_+}\) has \(\min_{\mathbf n}\Delta J_{\mathbf n}^2=0\), hence
	\[
	\xi_{S,\mathrm{odd}}^2(t)=0\quad\text{for all }t,\qquad
	\big\langle \xi_{S,\mathrm{odd}}^2\big\rangle_{\mathrm{ss}}=0.
	\]
	
	Finally, starting from \(\ket{++}\), the first (physical) click selects the parity sector with probability \(1/2\) each.
	Averaging the stationary values over sectors gives the overall stationary mean
	\[
	\boxed{\ \big\langle \xi_S^2\big\rangle_{\mathrm{ss,\,overall}}
		=\tfrac{1}{2}\Big(1-\tfrac{\eta}{2-\eta}\Big)+\tfrac{1}{2}\cdot 0
		=\frac{1-\eta}{2-\eta}\ }.
	\]
	
	
\section{Generalized CSS initial state (\texorpdfstring{$\phi=0$}{phi=0})}

We now take as initial state a two-spin coherent spin state (CSS) with polar angle \(\theta\) and \(\phi=0\):
\[
\ket{\psi(\theta)}=\ket{n}^{\otimes 2},\qquad 
\ket{n}=\sin\frac{\theta}{2}\ket{1}+\cos\frac{\theta}{2}\ket{0}.
\]
The jump operators remain
\[
\hat C_{\pm}=\sqrt{\gamma}\,\frac{\hat S_z^{(1)}\pm \hat S_z^{(2)}}{\sqrt{2}}
=\sqrt{\frac{\gamma}{8}}\big(\hat\sigma_z^{(1)}\pm \hat\sigma_z^{(2)}\big),
\qquad
\hat\Pi_z=\hat\sigma_z^{(1)}\otimes \hat\sigma_z^{(2)},
\]
with the convention \(\hat\sigma_z\ket{1}=+\ket{1}\), \(\hat\sigma_z\ket{0}=-\ket{0}\).

\subsection*{Jump probabilities from the post-jump norm}

Using \(\hat C_\pm^\dagger\hat C_\pm=\tfrac{\gamma}{4}(\hat{\mathbb I}\pm \hat\Pi_z)\),
\[
p_\pm(\theta)=\mel{\psi(\theta)}{\hat C_\pm^\dagger\hat C_\pm}{\psi(\theta)}
=\frac{\gamma}{4}\Big(1\pm \expval*{\hat\Pi_z}_{\psi(\theta)}\Big).
\]
Since \(\ket{\psi(\theta)}=\ket{n}\otimes\ket{n}\) with \(\expval*{\hat\sigma_z}_{\ket{n}}=-\cos\theta\),
\[
\expval*{\hat\Pi_z}_{\psi(\theta)}
=\expval*{\hat\sigma_z}_{\ket{n}}^2=\cos^2\theta.
\]
Hence
\[
\boxed{\,p_+(\theta)=\frac{\gamma}{4}\big(1+\cos^2\theta\big),\qquad
	p_-(\theta)=\frac{\gamma}{4}\sin^2\theta\,},\qquad
p_++p_-=\frac{\gamma}{2}.
\]

\subsection*{Normalized post-jump states}

Expanding \(\ket{\psi(\theta)}\) in the computational basis and acting with
\(\hat\sigma_z^{(1)}\pm \hat\sigma_z^{(2)}\), one finds (up to the overall factor \(\sqrt{\gamma/8}\))
\begin{align*}
	(\hat\sigma_z^{(1)}+\hat\sigma_z^{(2)})\ket{\psi(\theta)}
	&=2\sin^2\!\frac{\theta}{2}\ket{11}-2\cos^2\!\frac{\theta}{2}\ket{00},\\
	(\hat\sigma_z^{(1)}-\hat\sigma_z^{(2)})\ket{\psi(\theta)}
	&=2\sin\!\frac{\theta}{2}\cos\!\frac{\theta}{2}\,(\ket{10}-\ket{01}).
\end{align*}
Therefore
\[
\ket{\psi_+'}=
\frac{\sin^2(\tfrac{\theta}{2})\ket{11}-\cos^2(\tfrac{\theta}{2})\ket{00}}
{\sqrt{\sin^4(\tfrac{\theta}{2})+\cos^4(\tfrac{\theta}{2})}},
\qquad
\ket{\psi_-'}=\frac{\ket{10}-\ket{01}}{\sqrt{2}}=\ket{\psi_-}.
\]

\subsection*{Concurrence after one observed click and channel-averaged value}

For a pure state \(\ket{\Psi}=a\ket{00}+b\ket{01}+c\ket{10}+d\ket{11}\),
\(\mathcal C(\ket{\Psi})=2|ad-bc|\).
Thus
\[
\mathcal C_-(\theta)=1,\qquad
\mathcal C_+(\theta)=\frac{2\,\sin^2(\tfrac{\theta}{2})\cos^2(\tfrac{\theta}{2})}
{\sin^4(\tfrac{\theta}{2})+\cos^4(\tfrac{\theta}{2})}
=\frac{\sin^2\theta}{\,2-\sin^2\theta\,}.
\]
Averaging over the two channels with weights \(p_\pm/(p_++p_-)\) gives
\[
\boxed{\;
	\overline{\mathcal C}(\theta)=
	\frac{p_+\,\mathcal C_+ + p_-\,\mathcal C_-}{p_++p_-}
	=\frac{1}{2}\left[
	\sin^2\theta+\big(1+\cos^2\theta\big)\frac{\sin^2\theta}{\,2-\sin^2\theta\,}
	\right].\;}
\]
Checks: \(\overline{\mathcal C}(0)=\overline{\mathcal C}(\pi)=0\), \(\overline{\mathcal C}(\tfrac{\pi}{2})=1\).

\subsection*{Kitagawa--Ueda spin squeezing after one observed click}

For \(N=2\), \(\displaystyle \xi_S^2=\frac{4}{N}\min_{\mathbf n\perp \expval*{\hat{\mathbf J}}}\Delta J_{\mathbf n}^2\),
with \(\hat J_\alpha=\tfrac{1}{2}(\hat\sigma_\alpha^{(1)}+\hat\sigma_\alpha^{(2)})\).
In the \(-\) channel the post-jump state is the singlet \(\ket{\psi_-}\), for which
\(\Delta J_\alpha^2=0\) for all \(\alpha\) and \(\xi_{S,-}^2=0\).

In the \(+\) channel, for \(\ket{\psi_+'}=a\ket{00}+d\ket{11}\) (real \(a,d\)) one has
\(\expval*{\hat J_x}=\expval*{\hat J_y}=0\), \(\expval*{\hat J_z}=|d|^2-|a|^2\), and
\[
\Delta J_x^2=\frac{1}{2}\big(1+\expval*{\hat\sigma_x\otimes \hat\sigma_x}\big),\quad
\Delta J_y^2=\frac{1}{2}\big(1+\expval*{\hat\sigma_y\otimes \hat\sigma_y}\big).
\]
A short calculation yields
\[
\expval*{\hat\sigma_x\otimes \hat\sigma_x}=-\frac{\sin^2\theta}{\,2-\sin^2\theta\,},\qquad
\expval*{\hat\sigma_y\otimes \hat\sigma_y}=+\frac{\sin^2\theta}{\,2-\sin^2\theta\,},
\]
so that
\[
\Delta J_x^2=\frac{\cos^2\theta}{\,2-\sin^2\theta\,},\qquad
\Delta J_y^2=\frac{1}{\,2-\sin^2\theta\,}.
\]
Since \(\expval*{\hat{\mathbf J}}\) is along \(z\), the minimum transverse variance is \(\min(\Delta J_x^2,\Delta J_y^2)=\Delta J_x^2\).
With \(N=2\),
\[
\boxed{\,\xi_{S,+}^2(\theta)=2\,\Delta J_x^2
	=\frac{2\cos^2\theta}{\,2-\sin^2\theta\,}.}
\]
Therefore the channel-averaged KU parameter is
\[
\boxed{\;
	\overline{\xi_S^2}(\theta)
	=\frac{p_+\,\xi_{S,+}^2+p_-\,\xi_{S,-}^2}{p_++p_-}
	=\frac{1}{2}\big(1+\cos^2\theta\big)\,\frac{2\cos^2\theta}{\,2-\sin^2\theta\,}
	=\frac{(1+\cos^2\theta)\cos^2\theta}{\,2-\sin^2\theta\,}.}
\]
Edge cases: \(\overline{\xi_S^2}(0)=\overline{\xi_S^2}(\pi)=1\) (product states, no squeezing);
\(\overline{\xi_S^2}(\tfrac{\pi}{2})=0\) (both channels yield perfectly number-squeezed Bell states).

\medskip
In plots we set \(\gamma=1\) (dimensionless rates); the formulas above are independent of \(\gamma\) once probabilities are normalized by \(p_+{+}p_-=\gamma/2\).

	
\end{document}
